\documentclass[11pt,a4paper]{article}
\usepackage{fontspec}
\usepackage{polyglossia}
\setmainlanguage{polish}
\usepackage{geometry}
\usepackage{listings}
\usepackage{xcolor}
\usepackage{graphicx}
\usepackage{amsmath}
\usepackage{hyperref}
\usepackage{booktabs}
\usepackage{titlesec}

\geometry{margin=2cm}

\titleformat{\section}{\large\bfseries}{\thesection.}{0.5em}{}

\definecolor{codegreen}{rgb}{0,0.6,0}
\definecolor{codegray}{rgb}{0.5,0.5,0.5}
\definecolor{codepurple}{rgb}{0.58,0,0.82}
\definecolor{backcolour}{rgb}{0.95,0.95,0.92}

\lstset{
    language=Java,
    backgroundcolor=\color{backcolour},   
    commentstyle=\color{codegreen},
    keywordstyle=\color{magenta},
    numberstyle=\tiny\color{codegray},
    stringstyle=\color{codepurple},
    basicstyle=\ttfamily\small,
    breakatwhitespace=false,         
    breaklines=true,                 
    captionpos=b,                    
    keepspaces=true,                 
    numbers=left,                    
    numbersep=5pt,                  
    showspaces=false,                
    showstringspaces=false,
    showtabs=false,                  
    tabsize=2
}

\title{\textbf{Szczegółowe Sprawozdanie z Projektu Nr 1} \\ \Large Wielowątkowość wysokopoziomowa w języku Java}
\author{Analiza porównawcza technik klasycznych i wysokopoziomowych}
\date{\today}

\begin{document}

\maketitle

\tableofcontents
\newpage

\section{Cel i Zakres Projektu}
Głównym celem projektu było przeprowadzenie dogłębnej analizy porównawczej mechanizmów współbieżności w języku Java. Zbadano:
\begin{itemize}
    \item Podejście klasyczne: Bezpośrednie użycie interfejsu \texttt{Runnable} oraz klasy \texttt{Thread}.
    \item Podejście wysokopoziomowe: Wykorzystanie interfejsów \texttt{Executor}, \texttt{ExecutorService} oraz \texttt{ScheduledExecutorService} z pakietu \texttt{java.util.concurrent}.
\end{itemize}

\section{Weryfikacja Zgodności z Poleceniem}
Analiza przykładowego rozwiązania wykazuje pełną zgodność z wymaganiami określonymi w pliku \texttt{polecenie.txt}:
\begin{enumerate}
    \item \textbf{Zadanie programistyczne}: Zaimplementowano animację 100 elementów graficznych poruszających się po panelu Swing. Dodatkowo, każdy wątek wykonuje intensywne obliczenia matematyczne (\texttt{heavyCalculations}), co pozwala na realne obciążenie CPU.
    \item \textbf{Implementacja klasyczna}: Metoda \texttt{startThreads()} tworzy listę obiektów \texttt{Thread}, z których każdy zarządza jednym zadaniem.
    \item \textbf{Implementacja wysokopoziomowa}: Metody \texttt{startExecutorService()} oraz \texttt{startScheduled()} wykorzystują pule wątków (\texttt{ThreadPoolExecutor}) o stałym rozmiarze (domyślnie 4).
    \item \textbf{Analiza profilerem}: Przeprowadzono pomiary z wykorzystaniem narzędzia VisualVM 2.2.1 oraz narzędzi systemowych JDK (jstack, jstat).
\end{enumerate}

\section{Szczegóły Implementacji}

\subsection{Logika Zadania (KrokAnimacji)}
Każde zadanie wykonuje pętlę, w której zamiast pasywnego oczekiwania (\texttt{sleep}), wywoływana jest funkcja obciążająca procesor:
\begin{lstlisting}
public static void heavyCalculations() {
    double wynik = 0;
    for (int i = 0; i < 50000; i++) {
        wynik += Math.sin(i) * Math.tan(i);
    }
}
\end{lstlisting}
To podejście pozwala na badanie stanów wątków w warunkach wysokiej rywalizacji o zasoby.

\subsection{Pula Wątków (ExecutorService)}
W odróżnieniu od metody klasycznej, gdzie powstaje 100 wątków, w metodzie wysokopoziomowej zadania są kolejkowane do puli 4 wątków.
\begin{lstlisting}
executorService = Executors.newFixedThreadPool(ThreadingProject_P1.pula);
for (RuchomyElement el : panel.getElements()) {
    executorService.submit(new KrokAnimacji(el, panel, this));
}
\end{lstlisting}

\section{Analiza Profilerem (VisualVM)}

\subsection{Zużycie Procesora (CPU Usage)}
\begin{itemize}
    \item \textbf{Metoda Klasyczna}: Gwałtowny wzrost obciążenia CPU. System operacyjny musi zarządzać przełączaniem kontekstu dla 100 wątków konkurujących o rdzenie procesora. Widoczny narzut na zarządzanie wątkami.
    \item \textbf{ExecutorService (FixedThreadPool)}: Obciążenie stabilizuje się na poziomie odpowiadającym liczbie rdzeni/wątków w puli. Praca jest wykonywana płynnie, bez nadmiernego przełączania kontekstu.
\end{itemize}

\subsection{Stany Wątków (Thread States)}
Analiza zrzutu wątków (\texttt{jstack}) wykazała:
\begin{table}[h]
\centering
\begin{tabular}{@{}lll@{}}
\toprule
\textbf{Cecha} & \textbf{Klasycznie} & \textbf{Wysokopoziomowo} \\ \midrule
Liczba wątków robotników & 100 & 4 (stała) \\
Stan podczas pracy & RUNNABLE & RUNNABLE \\
Narzut pamięciowy & Wysoki (stack dla każdego wątku) & Niski (stała liczba stosów) \\
Zarządzanie kończeniem & Ręczne (interrupty w pętli) & Uproszczone (\texttt{shutdownNow()}) \\ \bottomrule
\end{tabular}
\caption{Porównanie charakterystyki wątków}
\end{table}

\section{Wnioski Końcowe}
1. \textbf{Wydajność}: Pule wątków są znacznie bardziej efektywne przy dużej liczbie drobnych zadań. Zapobiegają one zjawisku \textit{thread starvation} oraz nadmiernemu zużyciu pamięci RAM.
2. \textbf{Skalowalność}: API \texttt{java.util.concurrent} pozwala na łatwą zmianę strategii wykonywania zadań bez modyfikacji logiki samego zadania.
3. \textbf{Użyteczność}: \texttt{ScheduledExecutorService} jest jedynym rozsądnym wyborem dla zadań cyklicznych, zastępując błędogenne pętle z \texttt{Thread.sleep()}.

\end{document}
